% Part: first-order-logic
% Chapter: syntax-and-semantics
% Section: substitution

\documentclass[../../../include/open-logic-section]{subfiles}

\begin{document}

\olfileid{fol}{syn}{sub}

\olsection{প্রতিস্থাপন}

\begin{defn}[পদে প্রতিস্থাপন]
$s$-এ $x$-এর প্রতিটি সংঘটনের স্থলে $t$ \emph{প্রতিস্থাপন} করে পাওয়া
ফল $\Subst{s}{t}{x}$-কে আমরা পুনরাবৃত্তিমূলকভাবে সংজ্ঞায়িত করি:
\begin{enumerate}
\item \indcase{s}{c}{$\Subst{\indfrm}{t}{x}$ হলো শুধু $s$।}

\item \indcase{s}{y}{$y$ একটি চলরাশি এবং $y \not\ident x$ হলে
  $\Subst{\indfrm}{t}{x}$-ও শুধু~$s$।}

\item \indcase{s}{x}{$\Subst{\indfrm}{t}{x}$ হলো~$t$।}

\item\indcase{s}{\Atom{f}{t_1, \dots, t_n}}{$\Subst{\indfrmp}{t}{x}$ হলো
  $\Atom{f}{\Subst{t_1}{t}{x}, \dots, \Subst{t_n}{t}{x}}$।}
\end{enumerate}
\end{defn}

\begin{defn}
পদ~$t$ হলো $!A$-তে $x$-এর জন্য \emph{!!{free for}}, যদি $!A$-তে
$x$-এর কোনো মুক্ত সংঘটন এমন কোনো পরিমাণসূচকের পরিসরে না থাকে, যা
$t$-এর কোনো চলরাশিকে বদ্ধ করে।
\end{defn}

\begin{ex} ~
\begin{enumerate}
\item $\Obj v_8$, $\lexists[\Obj
  v_3]\Atom{\Obj A^2_4}{\Obj v_3,\Obj v_1}$-এ $\Obj v_1$-এর জন্য মুক্ত।

\item $\Obj f^2_1(\Obj v_1, \Obj v_2)$,
  $\lforall[\Obj v_2]\Atom{\Obj A^2_4}{\Obj v_0,\Obj v_2}$-এ
  $\Obj v_0$-এর জন্য \emph{মুক্ত নয়}।
\end{enumerate}
\end{ex}

\begin{defn}[!!a{formula}-এ প্রতিস্থাপন]
$!A$ যদি !!a{formula}, $x$ যদি !!a{variable} এবং $t$ যদি $!A$-তে
$x$-এর জন্য !!{free for} পদ হয়, তাহলে $\Subst{!A}{t}{x}$ হলো $!A$-তে
$x$-এর সব মুক্ত সংঘটনের স্থলে $t$ প্রতিস্থাপনের ফল।
\begin{enumerate}
\tagitem{prvFalse}{\indcase{!A}{\lfalse}{$\Subst{\indfrm}{t}{x}$ হলো
    $\lfalse$।}}{}

\tagitem{prvTrue}{\indcase{!A}{\ltrue}{$\Subst{\indfrm}{t}{x}$ হলো
    $\ltrue$।}}{}

\item \indcase{!A}{\Atom{P}{t_1,\dots,
    t_n}}{$\Subst{\indfrm}{t}{x}$ হলো $\Atom{P}{\Subst{t_1}{t}{x},
    \dots, \Subst{t_n}{t}{x}}$।}

\item \indcase{!A}{\eq[t_1][t_2]}{$\Subst{\indfrmp}{t}{x}$ হলো
  $\Subst{t_1}{t}{x} = \Subst{t_2}{t}{x}$।}

\tagitem{prvNot}{\indcase{!A}{\lnot !B}{$\Subst{\indfrmp}{t}{x}$ হলো
    $\lnot \Subst{!B}{t}{x}$।}}{}

\tagitem{prvAnd}{\indcase{!A}{(!B \land
    !C)}{$\Subst{\indfrmp}{t}{x}$ হলো $(\Subst{!B}{t}{x} \land
    \Subst{!C}{t}{x})$।}}{}

\tagitem{prvOr}{\indcase{!A}{(!B \lor
    !C)}{$\Subst{\indfrmp}{t}{x}$ হলো $(\Subst{!B}{t}{x} \lor
    \Subst{!C}{t}{x})$।}}{}

\tagitem{prvIf}{\indcase{!A}{(!B \lif
    !C)}{$\Subst{\indfrmp}{t}{x}$ হলো $(\Subst{!B}{t}{x} \lif
    \Subst{!C}{t}{x})$।}}{}

\tagitem{prvIff}{\indcase{!A}{(!B \liff
    !C)}{$\Subst{\indfrmp}{t}{x}$ হলো $(\Subst{!B}{t}{x} \liff
    \Subst{!C}{t}{x})$।}}{}

\tagitem{prvAll}{
\indcase{!A}{\lforall[y][!B]}{$y$ যদি $x$ ছাড়া অন্য চলরাশি হয়,
    তাহলে $\Subst{\indfrmp}{t}{x}$ হলো
    $\lforall[y][\Subst{!B}{t}{x}]$; অন্যথায়
    $\Subst{\indfrmp}{t}{x}$ শুধু $\indfrm$।}}{}

\tagitem{prvEx}{
\indcase{!A}{\lexists[y][!B]}{$y$ যদি $x$ ছাড়া অন্য চলরাশি হয়,
    তাহলে $\Subst{\indfrmp}{t}{x}$ হলো
    $\lexists[y][\Subst{!B}{t}{x}]$; অন্যথায়
    $\Subst{\indfrmp}{t}{x}$ শুধু $\indfrm$।}}{}
\end{enumerate}
\end{defn}

\begin{explain}
লক্ষ করুন, প্রতিস্থাপন নিষ্ফল হতে পারে: $x$ যদি $!A$-তে একেবারেই
সংঘটিত না হয়, তাহলে $\Subst{!A}{t}{x}$ হলো শুধু~$!A$।

$t$-কে $!A$-তে $x$-এর জন্য !!{free for} হওয়ার যে বিধিনিষেধ দেওয়া
হয়েছে, নিচের মতো ক্ষেত্র বাদ দেওয়ার জন্য তা প্রয়োজনীয়। যদি
$!A \ident \lexists[y][x < y]$ এবং $t \ident y$ হয়, তাহলে
$\Subst{!A}{t}{x}$ হতো $\lexists[y][y < y]$। এই ক্ষেত্রে মুক্ত
চলরাশি $y$ প্রতিস্থাপনের ফলে পরিমাণসূচক $\lexists[y]$-এর দ্বারা
“আবদ্ধ” হয়ে যায়, যা কাম্য নয়। যেমন, আমরা চাই যে
$\lforall[x][!B]$ সত্য হলেই $\Subst{!B}{t}{x}$-ও সত্য হবে। কিন্তু
$\lforall[x][\lexists[y][x < y]]$ বিবেচনা করুন (এখানে $!B$ হলো
$\lexists[y][x < y]$)। যেমন স্বাভাবিক সংখ্যার ক্ষেত্রে বাক্যটি সত্য:
প্রত্যেক সংখ্যা~$x$-এর চেয়ে বড় কোনো সংখ্যা~$y$ আছে। $x$-এর সম্ভাব্য
প্রতিস্থাপক হিসেবে $y$-কে অনুমোদন করলে আমরা
$\Subst{!B}{y}{x} \ident \lexists[y][y < y]$ পেতাম, যা মিথ্যা।
$t$-এর কোনো মুক্ত চলরাশি যেন $!A$-এর কোনো পরিমাণসূচকের দ্বারা শেষে
বদ্ধ না হয়ে যায়—এই শর্ত দিয়ে আমরা তা প্রতিরোধ করি।
\end{explain}

দুর্বহ সংকেতরীতি এড়াতে আমরা প্রায়ই নিচের প্রথা ব্যবহার করি। $!A$ যদি
এমন !!a{formula} হয় যাতে !!{variable}~$x$ মুক্তভাবে থাকতে পারে, তাহলে
তা বোঝাতে~$!A(x)$-ও লিখি। কোন $!A$ ও~$x$ বোঝানো হচ্ছে তা স্পষ্ট হলে,
এবং $t$ যদি একটি পদ হয় ($!A(x)$-এ $x$-এর জন্য মুক্ত ধরে), তাহলে
$\Subst{!A}{t}{x}$-এর সংক্ষিপ্ত রূপ হিসেবে $!A(t)$ লিখি। তাই যেমন আমরা
বলতে পারি, “$!A(t)$-কে $\lforall[x][!A(x)]$-এর একটি নিদর্শন বলি।”
এর অর্থ হলো, $!A$ যদি যেকোনো !!{formula}, $x$ যদি !!a{variable} এবং
$t$ যদি $!A$-তে $x$-এর জন্য মুক্ত পদ হয়, তাহলে $\Subst{!A}{t}{x}$
হলো $\lforall[x][!A]$-এর একটি নিদর্শন।

\end{document}
